\subsection{\classpypackageb}\label{sec:classpypackage}
The primary purpose of \gretl function packages is to extend
\gretl's repertoire of built-in procedures by adding new
estimators, hypothesis tests, or other analytical tools.
Although function packages may also serve other purposes
(e.g. pedagogy or replication exercises), those made available
via the \gretl server are expected to enhance \gretl's
functionality in non-trivial ways. More comprehensive information
is available at
\url{https://gretl.sourceforge.net/function_packages.html}.
For information on how to use function packages with \gretlpy,
please refer to Section~\ref{sec:gretl_contributed_packages}.


\subsubsection{\classpypackageb methods}
The \classpypackage class provides the following public methods:

\gretlfpycmd{get\_name()}
\estdesc{Returns package name.}
\begin{lstlisting}[style=pythonstylefoot]
.get_name() -> str
\end{lstlisting}

\gretlfpycmd{get\_functions()}
\estdesc{Returns the names of all public functions provided by the package.}
\begin{lstlisting}[style=pythonstylefoot]
.get_functions() -> list
\end{lstlisting}
\generaldesc{Note: all these functions are automatically bound
	to the user space once the \classpypackage\ instance is created.}

%\gretlfpycmd{set\_main()}
%\estdesc{XXX}
%\begin{lstlisting}[style=pythonstyle]
%.set_main(name=str)
%\end{lstlisting}
%\generaldesc{XXX}

\gretlfpycmd{show\_functions()}
\estdesc{Prints the signatures of all public functions provided by the package.}
\begin{lstlisting}[style=pythonstylefoot]
.show_functions()
\end{lstlisting}
\generaldesc{Using this method, the user may inspect
		a \gretl package, that is, retrieve detailed
		information about the functions it provides,
		including their names, arguments (with default
		values), and return types.}
